\documentclass[12pt,a4paper]{article}
\usepackage[T1]{fontenc}
\usepackage[utf8]{inputenc}
\usepackage[french]{babel}
\usepackage{geometry}
\usepackage{enumitem}
\usepackage{hyperref}
\usepackage{amsmath}
\usepackage[htt]{hyphenat}
\geometry{margin=2.5cm}
\setlist[itemize]{leftmargin=1.2cm}
\setlist[enumerate]{leftmargin=1.2cm}
\emergencystretch=4em
\tolerance=2000
\hbadness=4000
\sloppy

\title{Plan G\'en\'eral D\'etaill\'e du M\'emoire}
\author{}
\date{\today}

\begin{document}
\maketitle

\section{Contexte et positionnement du manuscrit}

Le pr\'esent plan d\'efinit l'orientation g\'en\'erale d'un m\'emoire consacr\'e au d\'eveloppement de m\'ethodes d'apprentissage automatique pour la d\'esagr\'egation temporelle et la mise \`a l'\'echelle spatiale des donn\'ees climatiques destin\'ees aux communaut\'es locales.

Le projet porte sur \textbf{Oracle}, une architecture causale et g\'en\'erative \`a deux \'etages pour le downscaling climatique probabiliste. Les axes structurants du manuscrit sont les suivants :
\begin{itemize}
    \item produire des champs climatiques haute r\'esolution \`a partir de pr\'edicteurs basse r\'esolution ;
    \item am\'eliorer la robustesse hors distribution dans un contexte de changement climatique ;
    \item limiter les hallucinations statistiques des mod\`eles purement corr\'elatifs par une structure causale explicite ;
    \item combiner un pr\'edicteur causal de moyenne haute r\'esolution avec un d\'ecodeur de diffusion r\'esiduelle.
\end{itemize}

\section{Plan global propos\'e du m\'emoire}

\subsection*{Pages pr\'eliminaire}
\begin{itemize}
    \item Page de garde
    \item D\'edicace (optionnel)
    \item Remerciements
    \item R\'esum\'e (fran\c{c}ais) + mots-cl\'es
    \item Abstract (anglais) + keywords
    \item Liste des abr\'eviations/sigles
    \item Table des mati\`eres
    \item Liste des figures et des tableaux
\end{itemize}

\subsection*{Corps du m\'emoire}
\begin{enumerate}
    \item \textbf{Introduction g\'en\'erale}
    \begin{itemize}
        \item Contexte climatique, besoins locaux et limites des r\'esolutions disponibles
        \item Probl\'ematique de g\'en\'eralisation hors distribution
        \item Questions de recherche
        \item Objectifs scientifiques et techniques
        \item Hypoth\`eses de travail
        \item M\'ethodologie globale
        \item Contributions attendues
        \item Organisation du document
    \end{itemize}

    \item \textbf{Chapitre 1 -- Fondements du downscaling climatique}
    \begin{itemize}
        \item Donn\'ees climatiques basse et haute r\'esolution
        \item D\'esagr\'egation temporelle et mise \`a l'\'echelle spatiale
        \item Mod\'elisation probabiliste des champs climatiques
        \item M\'etriques d'\'evaluation climatiques et probabilistes
        \item Trilemme : stabilit\'e, r\'ealisme des extr\^emes, robustesse OOD
        \item Synth\`ese de positionnement pour la suite
    \end{itemize}

    \item \textbf{Chapitre 2 -- \'Etat de l'art du downscaling par apprentissage profond}
    \begin{itemize}
        \item Approches d\'eterministes et U-Net
        \item Mod\`eles g\'en\'eratifs adversariaux et Residual cGAN
        \item Diffusion r\'esiduelle et EDM
        \item Contraintes physiques, XAI et causalit\'e
        \item Limites des approches existantes et gap scientifique
    \end{itemize}

    \item \textbf{Chapitre 3 -- Proposition Oracle : architecture two-stage causale}
    \begin{itemize}
        \item Graphe atmosph\'erique h\'et\'erog\`ene
        \item Encodeur intelligible de variables
        \item RCN : SCM, GRU et DAGMA
        \item Stage 1 : pr\'ediction de la moyenne causale
        \item Stage 2 : diffusion r\'esiduelle EDM
        \item Dynamique d'entra\^inement et propri\'et\'e (O3)
    \end{itemize}

    \item \textbf{Chapitre 4 -- Mat\'eriels, m\'ethodes, exp\'erimentations et discussion}
    \begin{itemize}
        \item Environnement logiciel et mat\'eriel (Python, PyTorch, Hydra, AMP, \texttt{torch.compile})
        \item Donn\'ees ACCESS-CM2 historique, pr\'etraitements et s\'equences temporelles ($\texttt{seq\_len}=16$, $\texttt{stride}=4$)
        \item Topologie Oracle : graphe h\'et\'erog\`ene, encodeur intelligible, RCN, Stage~1 (\texttt{GraphToGridDecoder}), Stage~2 (UNet EDM)
        \item Optimisation two-stage : AdamW + cosine pour Stage~1, Adam + groupes s\'epar\'es pour Stage~2, pertes \texttt{huber+cosine} et focal
        \item Diffusion EDM (Karras) : sampler de Heun, sigma calibr\'e
        \item Protocole d'\'evaluation : MAE, MSE, corr\'elation, CRPS, FSS, RAPSD, F1@p95/p99
        \item Diagnostic causal et v\'erification de la propri\'et\'e (O3) (ablation $A_{\text{dag}}$ vs $\mathbf{0}$)
        \item R\'esultats observ\'es, discussion critique et menaces \`a la validit\'e
    \end{itemize}

    \item \textbf{Conclusion g\'en\'erale et perspectives}
    \begin{itemize}
        \item R\'ecapitulatif des contributions
        \item R\'eponse \`a la probl\'ematique
        \item Limites du travail
        \item Perspectives de recherche et d'application
    \end{itemize}
\end{enumerate}

\subsection*{Fin du document}
\begin{itemize}
    \item Bibliographie
    \item Annexes (questionnaires, d\'etails techniques, r\'esultats compl\'ementaires)
\end{itemize}

\section{Plan d\'etaill\'e de l'Introduction g\'en\'erale}
\textbf{Titre recommand\'e :} \emph{Introduction g\'en\'erale}.

\subsection*{Objectif de l'introduction}
Poser le probl\`eme de recherche : comment produire des simulations climatiques locales haute r\'esolution qui soient \`a la fois r\'ealistes, probabilistes, physiquement coh\'erentes et robustes face aux climats futurs hors distribution.

\subsection*{Structure d\'etaill\'ee propos\'ee}

\subsubsection*{I.1 Contexte g\'en\'eral}
\begin{itemize}
    \item Importance des donn\'ees climatiques locales pour l'agriculture, la gestion des risques, les infrastructures et l'adaptation communautaire.
    \item Limitation pratique des mod\`eles globaux et r\'egionaux : r\'esolution souvent insuffisante pour les d\'ecisions locales.
    \item R\^ole du downscaling : transformer des champs atmosph\'eriques grossiers en cartes plus fines et exploitables.
\end{itemize}

\subsubsection*{I.2 Motivation}
\begin{itemize}
    \item Les mod\`eles d\'eterministes produisent souvent des champs trop liss\'es et sous-estiment les extr\^emes.
    \item Les mod\`eles g\'en\'eratifs am\'eliorent le r\'ealisme, mais peuvent rester instables, sous-dispersifs ou non interpr\'etables.
    \item Le changement climatique rend fragile l'hypoth\`ese de stationnarit\'e : les corr\'elations du pass\'e ne suffisent plus.
\end{itemize}

\subsubsection*{I.3 Probl\'ematique}
\begin{itemize}
    \item Les approches profondes apprennent parfois des raccourcis statistiques qui fonctionnent sur les donn\'ees historiques mais se d\'egradent en climat futur.
    \item La diffusion r\'esiduelle stabilise la g\'en\'eration, mais n'impose pas nativement de raisonnement causal ou physique.
    \item La question centrale est de concevoir une architecture o\`u la causalit\'e n'est pas un simple diagnostic ajout\'e apr\`es coup, mais un chemin de calcul n\'ecessaire \`a la pr\'ediction.
\end{itemize}

\subsubsection*{I.4 Questions de recherche}
\begin{itemize}
    \item Q1 : Les architectures g\'en\'eratives existantes, notamment Residual cGAN et ResDiff, suffisent-elles \`a garantir une robustesse hors distribution ?
    \item Q2 : Un mod\`ele causal structurel int\'egr\'e dans une dynamique r\'ecurrente permet-il de mieux guider le downscaling climatique ?
    \item Q3 : Une architecture two-stage, combinant moyenne causale et diffusion r\'esiduelle, permet-elle de pr\'eserver les extr\^emes et la calibration probabiliste ?
\end{itemize}

\subsubsection*{I.5 Objectifs}
\begin{itemize}
    \item \textbf{Objectif g\'en\'eral :} concevoir et \'evaluer une architecture d'apprentissage profond causale et g\'en\'erative pour le downscaling climatique local.
    \item \textbf{Objectifs sp\'ecifiques :}
    \begin{enumerate}
        \item analyser les limites des approches d\'eterministes, adversariales et diffusives ;
        \item formaliser le probl\`eme de downscaling probabiliste sous contrainte causale ;
        \item proposer une architecture Oracle \`a deux \'etages ;
        \item \'evaluer la qualit\'e spatiale, probabiliste et extr\^eme des sorties produites.
    \end{enumerate}
\end{itemize}

\subsubsection*{I.6 Hypoth\`eses de travail}
\begin{itemize}
    \item H1 : un mod\`ele qui apprend des relations causales invariantes g\'en\'eralise mieux qu'un mod\`ele purement corr\'elatif en contexte OOD.
    \item H2 : la s\'eparation moyenne causale / r\'esidu stochastique stabilise l'apprentissage et am\'eliore la repr\'esentation des extr\^emes.
    \item H3 : la propri\'et\'e d'indispensabilit\'e causale architecturale (O3) limite le risque de graphe causal simplement d\'ecoratif.
\end{itemize}

\subsubsection*{I.7 M\'ethodologie g\'en\'erale}
\begin{itemize}
    \item Revue critique des m\'ethodes de downscaling par apprentissage profond.
    \item Formalisation du probl\`eme comme estimation d'une distribution conditionnelle haute r\'esolution.
    \item Conception d'un graphe atmosph\'erique h\'et\'erog\`ene, d'un RCN et d'un d\'ecodeur de diffusion r\'esiduelle.
    \item \'Evaluation sur donn\'ees ERA5 avec m\'etriques d'erreur, de calibration et d'\'ev\'enements extr\^emes.
\end{itemize}

\subsubsection*{I.8 Contributions attendues}
\begin{itemize}
    \item Une architecture causale two-stage pour le downscaling climatique probabiliste.
    \item Un RCN associant SCM, GRU et contrainte d'acyclicit\'e DAGMA.
    \item Une interface causale entre le pr\'edicteur de moyenne et le d\'ecodeur de diffusion r\'esiduelle.
    \item Un protocole exp\'erimental centr\'e sur la calibration, les extr\^emes et les ablations causales.
\end{itemize}

\subsubsection*{I.9 Organisation du manuscrit}
\begin{itemize}
    \item R\'esumer en un paragraphe la fonction de chaque chapitre.
    \item Assurer une transition naturelle vers le chapitre 1.
\end{itemize}

\subsection*{Livrables attendus pour l'introduction}
\begin{itemize}
    \item Une formulation claire du probl\`eme de downscaling OOD.
    \item Trois questions de recherche directement reli\'ees \`a l'architecture propos\'ee.
    \item Une liste courte de contributions scientifiques et techniques.
    \item Un plan du manuscrit qui relie les fondements, l'\'etat de l'art, Oracle et l'\'evaluation.
\end{itemize}

\section{Plan d\'etaill\'e du Chapitre 1}
\textbf{Titre recommand\'e :} \emph{Fondements du downscaling climatique et concepts cibles}.

\subsection*{Objectif du chapitre}
Installer les bases climatiques, statistiques et m\'ethodologiques n\'ecessaires \`a la compr\'ehension du downscaling probabiliste et du besoin de causalit\'e dans les mod\`eles profonds.

\subsection*{Structure d\'etaill\'ee propos\'ee}

\subsubsection*{1.1 Introduction du chapitre}
\begin{itemize}
    \item Situer le chapitre comme socle conceptuel du m\'emoire.
    \item Annoncer les notions : donn\'ees climatiques, r\'esolution, downscaling, incertitude et OOD.
    \item Pr\'eciser que la revue des architectures avanc\'ees est trait\'ee au chapitre 2.
\end{itemize}

\subsubsection*{1.2 Donn\'ees climatiques et niveaux de r\'esolution}
\begin{itemize}
    \item Pr\'esenter les mod\`eles globaux, les mod\`eles r\'egionaux, les r\'eanalyses et les observations.
    \item Expliquer la diff\'erence entre basse r\'esolution et haute r\'esolution spatiale.
    \item Introduire les champs atmosph\'eriques utiles : vent, humidit\'e, temp\'erature, pression, pr\'ecipitations.
    \item Montrer l'int\'er\^et d'une repr\'esentation 3D de l'atmosph\`ere par niveaux de pression.
\end{itemize}

\subsubsection*{1.3 D\'esagr\'egation temporelle et mise \`a l'\'echelle spatiale}
\begin{itemize}
    \item D\'efinir la d\'esagr\'egation temporelle et la mise \`a l'\'echelle spatiale.
    \item Distinguer downscaling statistique, dynamique et hybride.
    \item Expliquer le passage typique d'une grille basse r\'esolution \`a une carte locale haute r\'esolution.
    \item Pr\'eciser le cas d'\'etude : entr\'ees LR, cible HR et s\'equences temporelles.
\end{itemize}

\subsubsection*{1.4 Downscaling probabiliste}
\begin{itemize}
    \item Formaliser l'objectif : estimer $p(x_{\mathrm{HR}} \mid y_{\mathrm{LR}}, s_{\mathrm{HR}}, \mathcal{G})$.
    \item Expliquer pourquoi une seule pr\'ediction moyenne est insuffisante pour les pr\'ecipitations et les extr\^emes.
    \item Introduire les notions de r\'esidu, incertitude, ensemble et calibration probabiliste.
\end{itemize}

\subsubsection*{1.5 Changement climatique et robustesse hors distribution}
\begin{itemize}
    \item D\'efinir la notion de distribution d'entra\^inement et de d\'ecalage hors distribution.
    \item Expliquer la fragilit\'e des corr\'elations apprises \`a partir du pass\'e.
    \item Motiver la recherche de m\'ecanismes invariants, li\'es \`a la physique et \`a la causalit\'e.
\end{itemize}

\subsubsection*{1.6 M\'etriques d'\'evaluation}
\begin{itemize}
    \item M\'etriques ponctuelles : MAE, MSE, corr\'elation spatiale.
    \item M\'etriques probabilistes : CRPS, spread--skill, calibration.
    \item M\'etriques spatiales et spectrales : FSS, RAPSD.
    \item M\'etriques d'extr\^emes : seuils p95/p99, pr\'ecision, rappel, F1-score.
\end{itemize}

\subsubsection*{1.7 Le trilemme du downscaling profond}
\begin{itemize}
    \item Stabilit\'e d'entra\^inement : \'eviter les \'echecs num\'eriques et les comportements adversariaux instables.
    \item R\'ealisme : reproduire les structures fines, la variance et les \'ev\'enements extr\^emes.
    \item Robustesse OOD : rester fiable lorsque les conditions climatiques futures diff\`erent de l'historique.
\end{itemize}

\subsubsection*{1.8 Synth\`ese et transition}
\begin{itemize}
    \item R\'esumer les concepts n\'ecessaires \`a la lecture de l'\'etat de l'art.
    \item Montrer que le probl\`eme exige plus qu'une simple super-r\'esolution d'image.
    \item Introduire les familles de mod\`eles profonds discut\'ees au chapitre 2.
\end{itemize}

\subsection*{Livrables attendus pour un Chapitre 1 solide}
\begin{itemize}
    \item Un sch\'ema basse r\'esolution vers haute r\'esolution.
    \item Un glossaire des variables climatiques et des m\'etriques.
    \item Un encadr\'e sur le probl\`eme OOD en climat.
    \item Une synth\`ese du trilemme stabilit\'e--r\'ealisme--robustesse.
\end{itemize}

\subsection*{Orientation r\'edactionnelle}
Le chapitre 1 doit rester \textbf{p\'edagogique et s\'electif} : il ne couvre pas toute la climatologie, mais fournit les bases n\'ecessaires \`a la compr\'ehension du probl\`eme et des choix m\'ethodologiques.

\section{Plan d\'etaill\'e du Chapitre 2}
\textbf{Titre recommand\'e :} \emph{\'Etat de l'art : approches d\'eterministes, g\'en\'eratives et causales pour le downscaling}.

\subsection*{Objectif du chapitre}
R\'ealiser une revue critique des approches de downscaling par apprentissage profond afin de faire appara\^itre le besoin d'une architecture causale, probabiliste et robuste hors distribution.

\subsection*{Structure d\'etaill\'ee propos\'ee}

\subsubsection*{2.1 Introduction du chapitre}
\begin{itemize}
    \item Rappeler le trilemme identifi\'e au chapitre 1.
    \item Pr\'esenter les axes de comparaison : stabilit\'e, extr\^emes, incertitude, physique, explicabilit\'e et OOD.
\end{itemize}

\subsubsection*{2.2 Approches d\'eterministes de downscaling}
\begin{itemize}
    \item Pr\'esenter les r\'egressions statistiques, U-Net et architectures de super-r\'esolution.
    \item Montrer leur capacit\'e \`a capturer les grandes structures.
    \item Expliquer le probl\`eme du lissage et de la sous-estimation des extr\^emes.
\end{itemize}

\subsubsection*{2.3 Mod\`eles g\'en\'eratifs adversariaux}
\begin{itemize}
    \item Introduire le cGAN et le Residual cGAN appliqu\'es au downscaling.
    \item Discuter les forces : textures plus fines, meilleure reproduction visuelle.
    \item Discuter les limites : sensibilit\'e \`a $\lambda_{\mathrm{adv}}$, sous-dispersion, mode collapse, hallucinations.
\end{itemize}

\subsubsection*{2.4 Diffusion r\'esiduelle et EDM}
\begin{itemize}
    \item Pr\'esenter ResDiff : moyenne d\'eterministe puis diffusion sur le r\'esidu.
    \item Expliquer l'int\'er\^et d'apprendre un r\'esidu de plus faible variance.
    \item Introduire EDM, les solveurs de d\'ebruitage et la calibration probabiliste.
    \item Montrer les limites : absence native de contrainte causale et risque de raccourcis statistiques.
\end{itemize}

\subsubsection*{2.5 Contraintes physiques et coh\'erence inter-variables}
\begin{itemize}
    \item Exemples de violations : $T_{\min} > T_{\max}$, incoh\'erences entre humidit\'e, pression et pr\'ecipitations.
    \item Approches par p\'enalit\'es physiques ou param\'etrisation contrainte.
    \item Limites des contraintes ajout\'ees sans structure causale interne.
\end{itemize}

\subsubsection*{2.6 Explicabilit\'e, causalit\'e et robustesse OOD}
\begin{itemize}
    \item Cartes de saillance et diagnostics post-hoc.
    \item Causalit\'e structurelle : SCM, DAG, DAGMA, invariance causale.
    \item Limite du graphe d\'ecoratif : un graphe peut exister sans devenir n\'ecessaire \`a la pr\'ediction.
\end{itemize}

\subsubsection*{2.7 Synth\`ese des limites P1--P5}
\begin{itemize}
    \item P1 : instabilit\'e ou sensibilit\'e d'entra\^inement.
    \item P2 : mauvaise quantification de l'incertitude et sous-dispersion.
    \item P3 : incoh\'erences physiques ou inter-variables.
    \item P4 : hypoth\`ese de stationnarit\'e fragile.
    \item P5 : causalit\'e non indispensable, ou graphe d\'ecoratif.
\end{itemize}

\subsubsection*{2.8 Gap de recherche et transition}
\begin{itemize}
    \item Formuler le besoin : une architecture o\`u le chemin causal gouverne la moyenne haute r\'esolution.
    \item Justifier la combinaison d'un RCN causal et d'une diffusion r\'esiduelle.
    \item Introduire Oracle comme r\'eponse architecturale au gap.
\end{itemize}

\subsection*{Livrables attendus pour le chapitre 2}
\begin{itemize}
    \item Un tableau comparatif des familles de mod\`eles : U-Net, cGAN, ResDiff, EDM, mod\`eles causaux.
    \item Une synth\`ese P1--P5 des limites de la litt\'erature.
    \item Une justification claire de l'architecture Oracle.
\end{itemize}

\section{Plan d\'etaill\'e du Chapitre 3}
\textbf{Titre recommand\'e :} \emph{Proposition Oracle : architecture two-stage causale pour le downscaling}.

\subsection*{Objectif du chapitre}
Pr\'esenter la contribution centrale du m\'emoire : une architecture o\`u un chemin causal construit une moyenne haute r\'esolution indispensable, puis une diffusion EDM mod\'elise uniquement le r\'esidu stochastique.

\subsection*{Structure d\'etaill\'ee propos\'ee}

\subsubsection*{3.1 Introduction du chapitre}
\begin{itemize}
    \item Rappeler les limites identifi\'ees dans l'\'etat de l'art.
    \item Introduire le principe two-stage : moyenne causale puis r\'esidu diffusif.
    \item Pr\'esenter la propri\'et\'e (O3) d'indispensabilit\'e causale architecturale.
\end{itemize}

\subsubsection*{3.2 Vue d'ensemble de l'architecture}
\begin{itemize}
    \item Pipeline global : graphe h\'et\'erog\`ene, encodeur intelligible, RCN, pr\'edicteur de moyenne, d\'ecodeur EDM.
    \item Entr\'ees basse r\'esolution, variables statiques haute r\'esolution et cible pr\'ecipitation.
    \item Sortie finale : $x_{\mathrm{HR}} = \mu_{\mathrm{HR}} + \delta$.
\end{itemize}

\subsubsection*{3.3 Graphe atmosph\'erique h\'et\'erog\`ene}
\begin{itemize}
    \item N\oe uds associ\'es aux niveaux de pression 850, 500 et 250 hPa.
    \item N\oe uds de surface et variables statiques haute r\'esolution.
    \item Ar\^etes spatiales, verticales, temporelles et statiques.
    \item Utilisation d'une structure compatible avec des donn\'ees climatiques multi-\'echelles.
\end{itemize}

\subsubsection*{3.4 Encodeur intelligible de variables}
\begin{itemize}
    \item Objectif : produire des repr\'esentations s\'emantiques lisibles par type de variable.
    \item M\'eta-chemins et agr\'egation de type SAGEConv.
    \item R\'eduction vers un ensemble restreint de variables causales interpr\'etables.
\end{itemize}

\subsubsection*{3.5 Recurrent Causal Network}
\begin{itemize}
    \item SCM : transition causale guid\'ee par une matrice d'adjacence acyclique.
    \item GRU : fusion entre plan causal et observation \`a l'instant courant.
    \item Perte de reconstruction et stabilisation temporelle de l'\'etat cach\'e.
    \item R\^ole du RCN : devenir le moteur logique de la simulation, et non un module d'explication externe.
\end{itemize}

\subsubsection*{3.6 DAGMA et projection spectrale}
\begin{itemize}
    \item Contrainte d'acyclicit\'e pour apprendre un DAG valide.
    \item P\'enalit\'e douce pendant l'entra\^inement.
    \item Projection dure pour rendre la structure causale certifiable.
\end{itemize}

\subsubsection*{3.7 Stage 1 : pr\'ediction de la moyenne causale}
\begin{itemize}
    \item Construction de $\mu_{\mathrm{HR}}$ \`a partir des \'etats causaux.
    \item Cross-attention vers une grille haute r\'esolution apprise.
    \item D\'ecodeur CNN Graph-to-Grid.
    \item Fonction de perte du stage 1 : erreur de moyenne, \'eventuelles contraintes physiques, r\'egularisation causale.
\end{itemize}

\subsubsection*{3.8 Interface inter-\'etages}
\begin{itemize}
    \item Calcul du r\'esidu initial $\delta_0 = x_{\mathrm{HR}} - \mu_{\mathrm{HR}}$.
    \item Recalibration de $\sigma_{\mathrm{data}}$ par estimation de type Welford.
    \item Gel du stage 1 pour emp\^echer la diffusion de court-circuiter le chemin causal.
    \item Gate d'ablation (O3) pour tester l'indispensabilit\'e de la causalit\'e.
\end{itemize}

\subsubsection*{3.9 Stage 2 : diffusion r\'esiduelle EDM}
\begin{itemize}
    \item Apprentissage du r\'esidu plut\^ot que du champ complet.
    \item Conditionnement par concat\'enation de canaux : r\'esidu bruit\'e, moyenne causale et baseline transform\'ee.
    \item Absence d'acc\`es direct \`a l'\'etat cach\'e causal dans le d\'ecodeur diffusif.
    \item Sampling EDM avec solveur de type Heun.
\end{itemize}

\subsubsection*{3.10 Pertes composites et dynamique d'entra\^inement}
\begin{itemize}
    \item Pertes causales, r\'ecurrentes, d\'eterministes, diffusives, physiques et de parcimonie.
    \item Entra\^inement s\'equentiel du stage 1 puis du stage 2.
    \item Gestion des gradients, clipping, pr\'ecision mixte et stabilit\'e num\'erique.
\end{itemize}

\subsubsection*{3.11 Synth\`ese et transition}
\begin{itemize}
    \item R\'ecapituler l'architecture Oracle.
    \item Relier chaque composant \`a une limite de la litt\'erature.
    \item Introduire le protocole exp\'erimental du chapitre 4.
\end{itemize}

\subsection*{Livrables attendus pour le chapitre 3}
\begin{itemize}
    \item Un sch\'ema global de l'architecture Oracle.
    \item Un algorithme d'entra\^inement en deux \'etages.
    \item Une description math\'ematique de $\mu_{\mathrm{HR}}$, $\delta_0$ et du conditionnement EDM.
    \item Une explication claire de la propri\'et\'e (O3).
\end{itemize}

\section{Plan d\'etaill\'e du Chapitre 4}
\textbf{Titre recommand\'e :} \emph{Mat\'eriels, m\'ethodes, exp\'erimentations et discussion}.

\subsection*{Objectif du chapitre}
Documenter les donn\'ees, l'environnement exp\'erimental, les hyperparam\`etres et les r\'esultats utilis\'es pour \'evaluer Oracle.

\subsection*{Structure d\'etaill\'ee propos\'ee}

\subsubsection*{4.1 Introduction du chapitre}
\begin{itemize}
    \item Rappeler les hypoth\`eses test\'ees (downscaling probabiliste, robustesse OOD, indispensabilit\'e causale (O3)).
    \item Annoncer les dimensions d'\'evaluation : erreur ponctuelle, calibration probabiliste, structure spectrale, extr\^emes (p95/p99) et ablations causales.
    \item Pr\'eciser que les choix exp\'erimentaux refl\`etent la configuration r\'eelle du d\'ep\^ot, pilot\'ee par Hydra (\texttt{config/training\_config.yaml}) et le point d'entr\'ee \texttt{ops/train\_st\_cdgm.py}.
\end{itemize}

\subsubsection*{4.2 Environnement logiciel et mat\'eriel}
\begin{itemize}
    \item Plateforme cible de r\'ef\'erence : Colab Pro avec GPU T4 (16 Go VRAM, 4--8 vCPU, 15--50 Go RAM) ; profil A100 80~Go support\'e via surcharge Hydra.
    \item Repli CPU : option d\'edi\'ee (\texttt{training.device=cpu, training.use\_amp=false, training.compile.enabled=false}).
    \item Pr\'ecision : entra\^inement en \emph{fp16} via AMP (\texttt{autocast} + \texttt{GradScaler}) sur GPU.
    \item Optimisation runtime : \texttt{torch.compile} activ\'e en mode \emph{default} pour encodeur, RCN et UNet ; \texttt{cudnn.benchmark=True} ; TF32 lorsque support\'e.
    \item Pile logicielle : Python 3.10, PyTorch $\geq 2.0$, PyTorch Geometric $\geq 2.3$, \texttt{diffusers} $\geq 0.21$, \texttt{accelerate}, \texttt{transformers}.
    \item Donn\'ees et calcul : NetCDF (\texttt{xarray}, \texttt{netCDF4}, \texttt{h5netcdf}), \texttt{dask} pour le chargement paresseux et \texttt{xbatcher} pour les batches spatiaux ; orchestration via \texttt{hydra-core} et \texttt{omegaconf}.
    \item Reproductibilit\'e : configuration s\'erialis\'ee par Hydra et persistance des checkpoints (\texttt{epoch\_last.pth}, \texttt{epoch\_best.pth}) avec \texttt{persist\_every=1}.
\end{itemize}

\subsubsection*{4.3 Jeux de donn\'ees et pr\'eparation}
\begin{itemize}
    \item Source : champs climatiques au format NetCDF, partition train avec \texttt{predictor\_ACCESS-CM2\_hist.nc} (LR) et \texttt{pr\_ACCESS-CM2\_hist.nc} (HR), pr\'edicteurs statiques ERA5 \texttt{ERA5\_eval\_ccam\_12km}.
    \item Entr\'ees LR : 15 champs dynamiques sur grille $23 \times 26$ : composantes $u$, $v$, $w$ du vent, humidit\'e sp\'ecifique $q$ et temp\'erature $t$ aux niveaux pression 850, 500 et 250~hPa.
    \item Champs statiques : orographie (\texttt{orog}), \'el\'evation effective (\texttt{he}) et indice de v\'eg\'etation (\texttt{vegt}), inject\'es dans le n\oe ud \texttt{SP\_HR} du graphe h\'et\'erog\`ene.
    \item Cible HR : pr\'ecipitations \texttt{pr} sur grille $172 \times 179$.
    \item Pr\'etraitements : standardisation moyenne/\'ecart-type des pr\'edicteurs, transformation $\log(1+p)$ des pr\'ecipitations (avec $\delta=0{,}01$~mm), strat\'egie de remplissage des NaN \texttt{mean}.
    \item Baseline interm\'ediaire : strat\'egie \texttt{hr\_smoothing} avec facteur 4, fournissant un champ de r\'ef\'erence dont le r\'esidu HR est ensuite mod\'elis\'e.
    \item S\'equences temporelles : \texttt{seq\_len}=16 et \texttt{stride}=4, charg\'ees \emph{lazy} (\texttt{eager\_load\_datasets=false}) pour limiter la pression m\'emoire NetCDF.
    \item Pipeline orient\'e graphe : conversion grille LR vers n\oe uds via \texttt{HeteroGraphBuilder} avec couches GP850/GP500/GP250 (\texttt{include\_mid\_layer=true}).
\end{itemize}

\subsubsection*{4.4 Topologie du mod\`ele}
\begin{itemize}
    \item M\'eta-chemins de l'encodeur intelligible : $5$ chemins dynamiques (\texttt{GP850\_spat\_adj}, \texttt{GP850\_to\_GP500}, \texttt{GP500\_spat\_adj}, \texttt{GP500\_to\_GP250}, \texttt{GP250\_spat\_adj}) plus, lorsque les statiques sont disponibles, le chemin causal \texttt{SP\_HR}~$\rightarrow$~\texttt{GP850} (for\c{c}age orographique).
    \item Encodeur : \texttt{hidden\_dim}=128, \texttt{conditioning\_dim}=128, m\'ecanisme \texttt{causal\_conditioning} actif avec \texttt{num\_dag\_tokens}=2 inject\'es dans la cross-attention.
    \item RCN : \texttt{hidden\_dim}=128, \texttt{driver\_dim}=15 (align\'e sur les variables LR), \texttt{dropout}=0{,}1, BPTT avec \texttt{detach\_interval}=8.
    \item Initialisation de $A_{\mathrm{dag}}$ : amorcage depuis le prior physique 6$\times$6 (couplage tridiagonal vertical + for\c{c}age orographique), avec un bruit gaussien $\sigma=0{,}02$.
    \item Stage~1 (\texttt{GraphToGridDecoder}) : architecture \texttt{graph\_to\_grid\_xattn} avec $d_\text{model}=128$, $4$ t\^etes d'attention, grille interm\'ediaire $43\times 45$ et raffinement CNN $128 \rightarrow 64 \rightarrow 32 \rightarrow 1$.
    \item Stage~2 (d\'ecodeur de diffusion, configuration V4 \emph{CorrDiff Large}) : UNet \texttt{diffusers} \texttt{UNet2DConditionModel} \`a 5 niveaux $172\times 179 \rightarrow 86\times 90 \rightarrow 43\times 45 \rightarrow 21\times 22 \rightarrow 10\times 11$, \texttt{block\_out\_channels}=[128, 256, 256, 256, 256] (\textasciitilde 53 millions de param\`etres), \texttt{layers\_per\_block}=2, blocs \texttt{DownBlock2D} aux 4 premiers niveaux puis \texttt{CrossAttnDownBlock2D} au niveau le plus petit (sym\'etrie miroir en remont\'ee), \texttt{attention\_head\_dim}=8, \texttt{norm\_num\_groups}=32, \texttt{use\_gradient\_checkpointing=true} pour rester sous 80~Go VRAM A100.
    \item Conditionnement : \texttt{causal\_concat\_mode=true} ; l'entr\'ee UNet est la concat\'enation par canal $(\delta_\sigma, \mu_{\mathrm{HR}}, \log(1+b))$, garantissant que le seul chemin par lequel $A_{\mathrm{dag}}$ atteint le d\'ecodeur est le canal $\mu_{\mathrm{HR}}$ (propri\'et\'e (O3) architecturale).
    \item Trajectoire architecturale V1$\rightarrow$V4 : V1 \emph{CorrDiff Mini} ([64,128,128], \textasciitilde 12~M params, Heun, stride=2), V2 \emph{CorrDiff Normal} 4 niveaux ([128,256,256,256], \textasciitilde 50~M params, tail-loss $w_{p95},w_{p99}=8,25$), V3 ajout EMA + stride=1 + horizon 100 \'epoques, V4 mont\'ee \`a 5 niveaux et tail-loss bump\'ee $w_{p95},w_{p99}=12,40$ (cf.\ section 4.8).
\end{itemize}

\subsubsection*{4.5 Protocole d'optimisation et hyperparam\`etres}
\begin{itemize}
    \item Architecture two-stage activ\'ee (\texttt{two\_stage.enabled=true}).
    \item \textbf{Stage~1 -- pr\'edicteur causal de moyenne :}
    \begin{itemize}
        \item Optimiseur \texttt{AdamW}, \texttt{weight\_decay}=$10^{-4}$, planificateur \texttt{cosine} avec 1 \'epoque de warmup.
        \item Taux d'apprentissage initial : $3 \times 10^{-4}$.
        \item Pertes pond\'er\'ees : $\lambda_{\text{reg}}=1{,}0$ (MSE sur $\mu_{\mathrm{HR}}$), $\beta_{\text{rec}}=0{,}05$ (reconstruction RCN), $\gamma_{\text{dag}}^{\max}=0{,}10$ avec rampe sur 2 \'epoques (DAGMA), $\lambda_{\text{L1}}=10^{-3}$, $\lambda_{\text{prior}}=0{,}05$ (ancrage sur le prior physique en \texttt{sum}).
        \item Plafond \texttt{epochs\_max}=7, early stopping plateau MSE (\texttt{patience}=3).
        \item Garde-fous anti-collapse : grad-gate $A_{\text{dag}}=1{,}0$, projection plancher $\|A\|_F \geq 0{,}10$, abort si $\|A\|_F < 0{,}05$.
        \item Diagnostic RAPSD log\'e \`a chaque \'epoque pour d\'etecter le sur-lissage.
    \end{itemize}
    \item \textbf{Stage~2 -- d\'ecodeur de diffusion r\'esiduelle (configuration V4) :}
    \begin{itemize}
        \item Optimiseur \texttt{AdamW}, taux d'apprentissage $2 \times 10^{-4}$, gradient clipping \`a $1{,}0$, AMP \emph{bf16} (CUDA).
        \item Stage~1 gel\'e (\texttt{freeze\_stage1=true}) : encodeur, RCN, DAGMA et \texttt{GraphToGridDecoder} sans gradient.
        \item Plafond \texttt{epochs\_max}=100 (V3/V4) sur stride=1 ($\sim 1{,}4\times 10^{4}$ \'echantillons par \'epoque), soit 110 \'epoques au total avec Stage~1 (10 + 100).
        \item Mode \texttt{causal\_concat\_mode=true} : entr\'ee UNet = concat\'enation $(\delta_\sigma, \mu_{\mathrm{HR}}, \log(1+b))$.
        \item Recalibration de $\sigma_{\text{data}}$ post-Stage~1 (\texttt{sigma\_data\_recalibrate=true}, \texttt{sigma\_min\_scale\_factor}=0{,}02), suivie d'une recalibration \emph{post-Stage~2} (\texttt{sigma\_data\_recalibrate\_post=true}) qui sert de filet de s\'ecurit\'e : si une d\'erive sup\'erieure \`a 10\,\% est d\'etect\'ee, le sampler d'inf\'erence est r\'ealign\'e.
        \item \textbf{Exponential Moving Average (BS37)} : moyenne exponentielle des poids du d\'ecodeur diffusif avec d\'ecalage $\beta_{\text{EMA}}=0{,}9999$, persist\'ee dans chaque checkpoint (cl\'e \texttt{diffusion\_ema\_state\_dict}). Le ph\'enom\`ene de moyenne EMA, propos\'e par \citet*{Karras2022EDM} et reformul\'e par \citet*{Karras2024EDM2}, est l'unique changement le plus impactant entre V2 et V3 : l'inf\'erence finale utilise les poids EMA, pas les poids \emph{live}.
        \item \textbf{Perte contrastive DAG (BS35)} : \`a chaque $N_c=4$ batchs, la perte EDM est \'evalu\'ee deux fois, sur $\mu_{\mathrm{HR}}(A_{\mathrm{dag}})$ et sur $\mu_{\mathrm{HR}}(\mathbf{0})$, et un terme de marge $\lambda_c \max(0, m_c - (\mathcal{L}_{\text{stage2}}^{(0)} - \mathcal{L}_{\text{stage2}}))$ avec $\lambda_c=0{,}5$ et $m_c=0{,}02$ force le d\'ecodeur \`a effectivement exploiter le chemin causal. Un garde-fou dynamique d\'ecro\^it $\lambda_c$ si la sensibilit\'e DAG reste n\'egative deux \'epoques cons\'ecutives, et un m\'ecanisme \emph{auto-tune} substitue l'ablation \texttt{zero} par \texttt{permute} puis \texttt{random} si le signal contrastif s'\'evanouit apr\`es 15 \'epoques.
    \end{itemize}
    \item \textbf{Param\`etres de la diffusion EDM (Karras 2022) :} formulation EDM, $\sigma_{\text{data}}$ recalibr\'e dynamiquement, $\sigma_{\min}=0{,}002$, $\sigma_{\max}=80$, $\rho=7$, $P_{\text{mean}}=-1{,}2$, $P_{\text{std}}=1{,}2$. Sampler d'inf\'erence : DPM-Solver++ multi-step \citep*{Lu2022DPMSolver} avec churn stochastique $S_{\text{churn}}=40$, 32 pas de d\'ebruitage, classifier-free guidance \texttt{cfg\_scale}=1{,}5 (avec \texttt{conditioning\_dropout\_prob}=0{,}1 \`a l'entra\^inement). La combinaison \emph{DPM-Solver++ + S\_churn=40 + CFG=1{,}5} constitue le \emph{Tier 0} V2-eval+ et apporte un gain net en RAPSD et Spread/RMSE sans aucun re-entra\^inement.
    \item \textbf{Pertes communes :} reconstruction EDM L2 \`a perte pond\'er\'ee tail-weighted (multipliateurs $w_{p95}=12$, $w_{p99}=40$ aux seuils $\tau_{95}=15$ et $\tau_{99}=30$~mm/jour) pour for\c{c}er l'apprentissage des extr\^emes ; le terme contrastif DAG ci-dessus est ajout\'e pour les batchs qualifi\'es.
    \item \textbf{R\'eglages d'entra\^inement :} \texttt{batch\_size}=64 effectif, \texttt{num\_workers}=8, \texttt{pin\_memory=true}, \texttt{prefetch\_factor}=2, planificateur cosine, persistance du checkpoint (\texttt{epoch\_last.pth}, \texttt{epoch\_best.pth}, plus \texttt{diffusion\_ema\_state\_dict}) \`a chaque \'epoque pour s\'ecuriser une \'eventuelle d\'econnexion Colab.
\end{itemize}

\subsubsection*{4.6 Protocole d'\'evaluation}
\begin{itemize}
    \item Erreurs ponctuelles : MAE, MSE, corr\'elation spatiale.
    \item Calibration probabiliste : CRPS estim\'e par \'echantillonnage Monte-Carlo (\texttt{num\_samples}=8 par d\'efaut, jusqu'\`a 32 membres pour CRPS exact) et diagnostic \emph{spread--skill}.
    \item Structure spatiale et spectrale : FSS multi-seuils, RAPSD log\'e \`a chaque \'epoque.
    \item \'Ev\'enements extr\^emes : F1-score aux percentiles $p95$ et $p99$ (\texttt{compute\_f1\_extremes=true}).
    \item Sampler d'\'evaluation : \texttt{dpm\_solver++} avec \texttt{eval\_num\_steps}=32 et $S_{\text{churn}}=40$ ; les poids du d\'ecodeur utilis\'es \`a l'inf\'erence sont la copie EMA ($\beta=0{,}9999$), pas les poids \emph{live}.
    \item Seuils op\'erationnels d\'eclar\'es dans la configuration : \texttt{MSE\_MAX}=0{,}15, \texttt{MAE\_MAX}=0{,}20, \texttt{CORR\_MIN}=0{,}91.
    \item Sorties : visualisations sauvegard\'ees (\texttt{save\_visualizations=true}) dans \texttt{results/}.
\end{itemize}

\subsubsection*{4.7 Diagnostic causal et v\'erification de la propri\'et\'e (O3)}
\begin{itemize}
    \item \textbf{Diagnostic amont post-Stage~1.} Ablation causale automatis\'ee (\texttt{causal\_ablation.enabled=true}, \texttt{n\_samples}=100) sur le \emph{champ moyen} produit par Stage~1, avec calcul de la sensibilit\'e relative $\delta\mu = \|\mu_{\mathrm{HR}}(A_{\mathrm{dag}}) - \mu_{\mathrm{HR}}(\mathbf{0})\|_2 / \|\mu_{\mathrm{HR}}(A_{\mathrm{dag}})\|_2$.
    \item Trois r\'egimes report\'es par le diagnostic : $\delta\mu < 0{,}01$ (\textsc{DAG d\'ecoratif} -- abort de la \emph{run}), $0{,}01 \leq \delta\mu < 0{,}10$ (\textsc{signal causal modeste}), $\delta\mu \geq 0{,}10$ (\textsc{signal causal fort}). Le mod\`ele V2 mesure $\delta\mu \approx 0{,}18$, plac\'e fermement dans le r\'egime \textsc{fort}.
    \item \textbf{Signal contrastif pendant l'entra\^inement (BS35-CONTRASTIVE).} Sensibilit\'e DAG $\Delta\mathcal{L}_{\text{stage2}} = \mathcal{L}_{\text{stage2}}^{(0)} - \mathcal{L}_{\text{stage2}}$ log\'ee \`a chaque pas du Stage~2. Une valeur positive certifie que le d\'ecodeur diffusif exploite activement $\mu_{\mathrm{HR}}$ ; le garde-fou progressif d\'ecrit en 4.5 attenue la perte contrastive si $\Delta\mathcal{L} < 0$ pour pr\'evenir la d\'estabilisation.
    \item \textbf{Suite d'interventions \`a l'inf\'erence (BS35).} Quatre modes \'evalu\'es sur le \emph{m\^eme} fold de validation, $K=16$ membres d'ensemble par \'echantillon :
    \begin{itemize}
        \item \textsc{Causal} : pipeline standard (DAG appris).
        \item \textsc{Zero-DAG} : substitution $A_{\mathrm{dag}} \leftarrow \mathbf{0}$ \`a l'inf\'erence uniquement.
        \item \textsc{Random-DAG} : substitution par une matrice triangulaire sup\'erieure al\'eatoire \`a sparsit\'e identique.
        \item \textsc{Permute-DAG} : permutation al\'eatoire des lignes de $A_{\mathrm{dag}}$.
    \end{itemize}
    \item \textbf{Contrefactuel architectural (\textsc{Reg.\ baseline}).} Substitution de tout le pipeline causal (encodeur + RCN + DAGMA + \texttt{GraphToGridDecoder}) par un pr\'edicteur de moyenne iso-capacit\'e \texttt{RegressionMeanPredictor} (UNet \emph{CorrDiff-style} \`a budget param\'etrique appari\'e), entra\^in\'e pour le m\^eme nombre d'\'epoques. Si \textsc{Reg.\ baseline} sous-performe \textsc{Causal} \`a budget identique, l'architecture causale est empiriquement n\'ecessaire ; sinon le DAG est r\'eduit au statut d'auxiliaire interpr\'etatif.
    \item Crit\`ere de la propri\'et\'e (O3) : $r_{O3} = \mathrm{F1@p99}^{(\textsc{Causal})} - \mathrm{F1@p99}^{(\textsc{Zero-DAG})} \geq 0{,}05$ ; abort de la \emph{run} si l'\'epreuve amont \'echoue (\texttt{abort\_if\_fail=true}). La V2 mesure $r_{O3} \approx 0{,}13$.
    \item Suivi continu de $\|A_{\mathrm{dag}}\|_F$ et de la sparsit\'e induite par DAGMA + L1 + projection spectrale.
\end{itemize}

\subsubsection*{4.8 R\'esultats exp\'erimentaux : trajectoire V1$\rightarrow$V4}
\begin{itemize}
    \item \textbf{Cibles Oracle (\citep{Mardani2024ResDiff, Rampal2024cGAN, Mardani2025CorrDiff}) :} Pearson $0{,}896$, F1@$p95$ $0{,}837$, F1@$p99$ $0{,}754$. Ces valeurs sont \emph{l'objectif d\'epasser}, pas un r\'esultat de la pr\'esente \'etude ; elles sont d\'eriv\'ees de la litt\'erature existante (cGAN r\'esiduel et ResDiff) et fixent la barre que la trajectoire V1$\rightarrow$V4 cherche \`a franchir.
    \item \textbf{Strat\'egie de \emph{full retrain isol\'e} :} chaque variante V1, V2, V3, V4 dispose de son propre r\'epertoire de \texttt{checkpoint} (\texttt{ckpt/}, \texttt{ckpt\_v2\_corrdiff\_normal/}, \texttt{ckpt\_v3\_ema\_long/}, \texttt{ckpt\_v4\_corrdiff\_large/}) afin d'\'eviter tout \emph{state\_dict mismatch} entre architectures et de pr\'eserver chaque baseline pour la comparaison m\'emoire.
    \item Tableau de synth\`ese (valeurs en \emph{(en cours)} \`a renseigner depuis les fichiers \texttt{v2\_eval\_plus\_metrics.json}, \texttt{v3\_metrics.json} et \texttt{v4\_metrics.json} produits par la cellule \texttt{FINAL\_VALIDATION} du notebook) :

\begin{center}
\small
\begin{tabular}{lcccc}
\hline
Variante  & Pearson $\uparrow$ & F1@$p95$ $\uparrow$ & F1@$p99$ $\uparrow$ & RAPSD dist. $\downarrow$ \\
\hline
V1 (CorrDiff Mini)         & $0{,}36$       & $0{,}42$       & $0{,}0037$     & $812$ \\
V2 (CorrDiff Normal, 50\,ep) & $0{,}71$       & $0{,}42$       & $0{,}29$       & $364$ \\
V2 étendu (200\,ep, stride=2) & $0{,}82$       & $0{,}65$       & $0{,}53$       & $258$ \\
V2-eval+ (Tier 0, inf.)    & (post-run)     & (post-run)     & (post-run)     & (post-run) \\
V3 (Tier 0+1, EMA + 100\,ep) & (post-run Colab) & (post-run Colab) & (post-run Colab) & (post-run Colab) \\
V4 (Tier 0+1+2, CorrDiff L) & (post-run Colab) & (post-run Colab) & (post-run Colab) & (post-run Colab) \\
\hline
Cible Oracle              & $0{,}896$      & $0{,}837$      & $0{,}754$      & $\rightarrow 0$ \\
\hline
\end{tabular}
\end{center}
    \item \textbf{Lecture de la trajectoire V1$\rightarrow$V2 (mesur\'ee).} Le passage de \emph{CorrDiff Mini} \`a \emph{CorrDiff Normal} avec ajout de la \emph{tail-loss} pond\'er\'ee multiplie la corr\'elation par $\sim 2$ et le F1@$p99$ par $\sim 78$, ce qui confirme que la sous-performance de V1 \'etait dict\'ee par la \emph{capacit\'e} de Stage~2 et non par le pr\'edicteur causal Stage~1.
    \item \textbf{Apport attendu de chaque tiers (V2$\rightarrow$V4).} Tier 0 (V2-eval+, gain \emph{gratuit} : DPM-Solver++ + $S_{\text{churn}}$ + CFG + EMA \`a l'inf\'erence) : $+0{,}07$ \`a $+0{,}11$ Pearson, $+0{,}10$ \`a $+0{,}20$ F1@$p99$. Tier 1 (V3, EMA pendant l'entra\^inement + stride=1 + horizon $100$ \'epoques + recalibration post-S2) : $+0{,}06$ \`a $+0{,}08$ Pearson, $+0{,}15$ \`a $+0{,}25$ F1@$p99$. Tier 2 (V4, mont\'ee \`a 5 niveaux + \emph{tail-loss} bump\'ee) : $+0{,}04$ \`a $+0{,}06$ Pearson, $+0{,}10$ \`a $+0{,}15$ F1@$p99$, gain principal sur RAPSD via la profondeur du \emph{champ r\'eceptif}.
    \item Les valeurs ci-dessus sont issues du notebook \texttt{st\_cdgm\_training\_evaluation.ipynb} et persist\'ees \`a chaque \'epoque ; les ex\'ecutions futures seront r\'eestim\'ees avec intervalles de confiance et r\'ep\'etitions \texttt{seed} multiples.
\end{itemize}

\subsubsection*{4.8b Validation causale (apport mesurable du DAG)}
\begin{itemize}
    \item \textbf{Apport scientifique central :} l'objectif n'est pas seulement de d\'epasser les m\'etriques Oracle, mais de \emph{quantifier l'apport du DAG} dans cette d\'epassement. L'argument repose sur quatre indicateurs cumulatifs :
    \begin{enumerate}
        \item Diagnostic amont $\delta\mu$ (cf.\ 4.7) -- mesure \emph{intrins\`eque} \`a Stage~1.
        \item Sensibilit\'e contrastive $\Delta\mathcal{L}_{\text{stage2}}$ (cf.\ 4.5/4.7) -- mesure \emph{dynamique} \`a l'entra\^inement de Stage~2.
        \item Suite BS35 d'interventions \`a l'inf\'erence (cf.\ 4.7) -- mesure \emph{descendante} sur la qualit\'e finale.
        \item Contrefactuel architectural \textsc{Reg.\ baseline} -- mesure \emph{contrefactuelle} \`a iso-capacit\'e.
    \end{enumerate}
    \item Tableau BS35 attendu (\`a renseigner depuis \texttt{v4\_bs35\_ablation.json}) :

\begin{center}
\small
\begin{tabular}{lccc}
\hline
Mode                       & Pearson $\uparrow$ & F1@$p99$ $\uparrow$ & $r_{O3}$ \\
\hline
\textsc{Causal} (V4)       & (en cours) & (en cours) & --- \\
\textsc{Zero-DAG} (V4)     & (en cours) & (en cours) & (en cours) \\
\textsc{Random-DAG} (V4)   & (en cours) & (en cours) & (en cours) \\
\textsc{Permute-DAG} (V4)  & (en cours) & (en cours) & (en cours) \\
\textsc{Reg.\ baseline}    & (en cours) & (en cours) & --- \\
\hline
\end{tabular}
\end{center}
    \item Une d\'efaillance de \textsc{Reg.\ baseline} \`a \'egaler \textsc{Causal} \`a budget appari\'e constituerait la \emph{plus forte} preuve empirique disponible que la causalit\'e est \emph{n\'ecessaire}, par opposition \`a un graphe d\'ecoratif. Cette ligne joue un r\^ole analogue au bras de contr\^ole d'une exp\'erience randomis\'ee.
\end{itemize}

\subsubsection*{4.9 Analyse des r\'esultats}
\begin{itemize}
    \item Lire la corr\'elation spatiale comme un indicateur de localisation des structures pluvieuses produites.
    \item Discuter le CRPS comme mesure conjointe de calibration et de finesse de l'ensemble (avec \'ecart \emph{spread--skill}).
    \item Justifier l'apport du focal loss et de la pond\'eration des extr\^emes pour les scores $F1@p95$ et $F1@p99$.
    \item Relier les r\'esultats au r\^ole du RCN (m\'emoire causale), du Stage~1 (champ moyen explicite) et du d\'ecodeur EDM (variabilit\'e fine).
    \item Discuter le compromis temps d'inf\'erence vs nombre de pas (\texttt{eval\_num\_steps}) et nombre de membres CRPS.
\end{itemize}

\subsubsection*{4.10 Discussion critique}
\begin{itemize}
    \item Risque de \emph{collapse} du r\'esidu si $\mu_{\mathrm{HR}}$ absorbe une fraction trop importante du signal HR ; le d\'eclenchement du garde-fou \texttt{abort\_on\_collapse} et de \texttt{dag\_floor\_projection} doit \^etre report\'e.
    \item Limite d'une baseline lissage HR pour certains r\'egimes convectifs ; \'evolution potentielle vers des baselines bicubiques calibr\'ees ou apprises.
    \item Difficult\'e des queues lourdes de pr\'ecipitations et int\'er\^et de variantes $t$-EDM ou WassDiff comme perspectives.
    \item Co\^ut d'inf\'erence dominant : compromis vitesse/pr\'ecision \`a r\'egler via \texttt{eval\_num\_steps}, \texttt{num\_samples} et \texttt{eval\_scheduler}.
    \item Sensibilit\'e des m\'etriques au choix du seuil de pr\'ecipitations et \`a la transformation $\log(1+p)$.
\end{itemize}

\subsubsection*{4.11 Menaces \`a la validit\'e}
\begin{itemize}
    \item Validit\'e interne : qualit\'e des splits temporels, choix de normalisation, traitement des NaN (strat\'egie \texttt{mean}), reproductibilit\'e des graines.
    \item Validit\'e externe : couverture limit\'ee aux donn\'ees ACCESS-CM2 historique, n\'ecessit\'e d'\'evaluations sur d'autres r\'eanalyses (ERA5) et sur les sc\'enarios CMIP6/CORDEX.
    \item Validit\'e statistique : incertitude des m\'etriques mono-run, n\'ecessit\'e de r\'ep\'etitions et d'intervalles bootstrap.
    \item Validit\'e op\'erationnelle : passage de l'apprentissage en climat historique au d\'eploiement sur sc\'enarios futurs (OOD), avec recalibration potentielle de $\sigma_{\text{data}}$ et v\'erification de stabilit\'e du DAG.
\end{itemize}

\subsubsection*{4.12 Synth\`ese du chapitre}
\begin{itemize}
    \item R\'esumer ce que les r\'esultats confirment vis-\`a-vis des hypoth\`eses Q1--Q3.
    \item Identifier les points encore \`a consolider (r\'ep\'etitions, OOD, baselines suppl\'ementaires).
    \item Pr\'eparer la transition vers la conclusion g\'en\'erale.
\end{itemize}

\subsection*{Livrables attendus pour le chapitre 4}
\begin{itemize}
    \item Tableau exhaustif des donn\'ees, des variables et des hyperparam\`etres effectivement utilis\'es (issu de \texttt{config/training\_config.yaml}).
    \item Tableau de r\'esultats principaux (point estimate + r\'ep\'etitions \`a venir).
    \item Figures comparatives : champs HR pr\'edits, observations, baseline, r\'esidu et $\mu_{\mathrm{HR}}$.
    \item Tableau d'ablations causales : DAG appris, DAG nul, DAG al\'eatoire, prior physique seul.
    \item Section de discussion critique, limites et menaces \`a la validit\'e.
\end{itemize}

\section{Plan d\'etaill\'e de la Conclusion g\'en\'erale}
\textbf{Titre recommand\'e :} \emph{Conclusion g\'en\'erale et perspectives}.

\subsection*{Objectif de la conclusion}
Fermer l'argument scientifique : revenir sur la probl\'ematique, synth\'etiser les apports valid\'es, expliciter les limites et ouvrir des perspectives r\'ealistes.

\subsection*{Structure d\'etaill\'ee propos\'ee}

\subsubsection*{C.1 R\'ecapitulatif de la probl\'ematique}
\begin{itemize}
    \item Reprendre le probl\`eme : produire un downscaling local probabiliste fiable en climat futur.
    \item Rappeler le fil directeur : du trilemme OOD \`a l'architecture Oracle.
\end{itemize}

\subsubsection*{C.2 Synth\`ese des contributions}
\begin{itemize}
    \item Contribution conceptuelle : formulation du besoin de causalit\'e indispensable, au-del\`a de l'explicabilit\'e post-hoc.
    \item Contribution architecturale : graphe h\'et\'erog\`ene, encodeur intelligible, RCN, DAGMA et projection spectrale.
    \item Contribution g\'en\'erative : d\'ecomposition moyenne causale / r\'esidu EDM.
    \item Contribution exp\'erimentale : protocole centr\'e sur calibration, extr\^emes et ablations causales.
\end{itemize}

\subsubsection*{C.3 R\'eponses aux questions de recherche}
\begin{itemize}
    \item R\'eponse \`a Q1 : positionner les limites des mod\`eles adversariaux et diffusifs non causaux.
    \item R\'eponse \`a Q2 : expliquer le r\^ole du RCN comme moteur causal r\'ecurrent.
    \item R\'eponse \`a Q3 : relier les r\'esultats aux m\'etriques de calibration et d'extr\^emes.
\end{itemize}

\subsubsection*{C.4 Limites du travail}
\begin{itemize}
    \item Domaine climatique et volume de donn\'ees encore limit\'es.
    \item Validation OOD \`a consolider avec des sc\'enarios climatiques futurs.
    \item Baseline non causale \`a am\'eliorer pour certains usages op\'erationnels.
    \item Co\^ut d'inf\'erence de la diffusion \`a optimiser.
\end{itemize}

\subsubsection*{C.5 Perspectives}
\begin{itemize}
    \item Extension multivariable avec pertes physiques inter-variables.
    \item \'Evaluation sur plusieurs domaines climatiques et jeux CORDEX/CMIP6 (g\'en\'eralisation OOD r\'eelle au sens de la section 1.5).
    \item \textbf{Tier 3 architectural :} entra\^inement \emph{joint} de Stage~1 et Stage~2 (sans gel apr\`es Stage~1) sous contrainte explicite de la propri\'et\'e (O3) ; ajout d'une perte \emph{spectrale} sur le RAPSD pour pousser au plus pr\`es la cible Oracle ; passage \`a $t$-EDM \citep*{Pandey2024tEDM} pour les queues lourdes Student-$t$, ou \`a WassDiff \citep*{Liu2024WassDiff} pour la calibration multivari\'ee.
    \item Couplage avec correction de biais et cha\^ines d'impact pour les communaut\'es locales.
    \item Outils de visualisation XAI pour auditer les chemins causaux appris (carte de chaleur de la sensibilit\'e \`a chaque ar\^ete de $A_{\mathrm{dag}}$, projection saillance-DAG sur la moyenne $\mu_{\mathrm{HR}}$).
\end{itemize}

\subsubsection*{C.6 Mot de cl\^oture}
\begin{itemize}
    \item Rappeler la valeur scientifique : passer d'un mod\`ele qui reproduit des motifs climatiques \`a un mod\`ele qui exploite des m\'ecanismes causaux invariants.
    \item Rappeler la valeur pratique : produire des cartes locales plus fiables pour l'adaptation climatique.
\end{itemize}

\subsection*{Livrables attendus pour la conclusion}
\begin{itemize}
    \item Un tableau ``Questions de recherche vs r\'eponses apport\'ees''.
    \item Une section limites courte, explicite et honn\^ete.
    \item Une liste de perspectives prioris\'ees pour poursuite en article ou th\`ese.
\end{itemize}

\end{document}
